\section{Introduction}
\label{p1:intro}

A long constant run in a sequence of independent fair signs starts
with probability \(2^{-L}\). There are \(N\) possible starts in a
dyadic interval, so the scale \(L=\log_2 N+O(1)\) produces a bounded
expected count. Complete multiplicativity preserves this scale but
creates relations between distant windows. For example, the signs
at \(2u\) and \(3u\) inherit the same sign at \(u\). The question
addressed here is how these arithmetic relations first alter the
Poisson law of the total number of runs.

The leading deviation comes from a sparse family of pairs and has
size \(L N^{-2/3}\). Once the exact pair correction is retained,
the error has optimal power \(N^{-4/3}\). Its obstruction is the
product of two pair interactions. Triple interactions are smaller:
the third factorial cumulant has optimal power \(N^{-41/30}\).
Identifying a rational pair coefficient is a two-window calculation.
Controlling the complete law also requires estimates uniform in
the number of windows, and arithmetic restrictions on how much
information three windows can share.

\subsection{Model and main results}
\label{p1:intro:model}

Choose independent random variables \(f(\ell)\), indexed by primes
\(\ell\), with
\[
 \P(f(\ell)=1)=\P(f(\ell)=-1)=\tfrac12,
 \qquad
 f(n)=\prod_{\ell}\! f(\ell)^{v_\ell(n)}\quad(n\ge1).
\]
Thus \(f\) is completely multiplicative and \(f(n^2)=1\).
Fix \(C>0\). Throughout the paper \(N\ge2\) and \(L\ge1\) are integers,
\begin{equation}
 I_N=[N,2N)\cap\Z,\qquad
 |L-\log_2 N|\le C,\qquad B=L+1,\qquad
 p=2^{-L},\qquad \lambda_N=Np.
 \label{p1:eq:critical}
\end{equation}
All logarithms without a subscript are natural. In particular,
\(\lambda_N\) stays in a compact subinterval of \((0,\infty)\);
it need not converge. Define
\begin{equation}
 J_x=\ind_{\{f(x-1)=-f(x),\ 
             f(x)=f(x+1)=\cdots=f(x+L-1)\}},
 \qquad Z_N=\sum_{x\in I_N}J_x .
 \label{p1:eq:model}
\end{equation}
The left boundary condition ensures that a run is counted once,
at its start. Both signs of the run are included.
The individual probabilities \(\E J_x\) are not assumed equal to \(p\).

For a finite signed measure \(\nu\) on \(\N=\{0,1,2,\ldots\}\), set
\(\|\nu\|_1=\sum_{k\ge0}|\nu\{k\}|\).
Our total variation convention for probabilities is
\(\TV(P,Q)=\frac12\|P-Q\|_1\).
Write
\[
 \Pi_\mu=\Pois(\mu),\qquad
 \Delta=\delta_1-\delta_0,\qquad
 \mu_N=\E Z_N,\qquad
 \alpha_N=\Var(Z_N)-\mu_N .
\]
Products of measures are convolutions, and
\(\Delta^{*0}=\delta_0\). For \(r\ge1\), the falling factorial is
\((z)_r=z(z-1)\cdots(z-r+1)\), with \((z)_0=1\).
The second-order reference is
\begin{equation}
 Q_{\mu,\alpha}
    =\Pi_\mu+\frac{\alpha}{2}\Pi_\mu*\Delta^{*2}.
 \label{p1:eq:Q}
\end{equation}
Its total mass is one. Its positivity in the range needed below
will follow from an exact criterion.

\begin{theorem}[Optimal quadratic Poisson--Charlier approximation]
\label{p1:thm:main}
Fix \(C>0\). For every \(\varepsilon>0\), there are
positive constants \(c_{C,\varepsilon}\), \(C_{C,\varepsilon}\)
such that, for all sufficiently large \(N\), uniformly over
\((N,L)\) satisfying \eqref{p1:eq:critical},
\begin{equation}
 c_{C,\varepsilon}N^{-4/3-\varepsilon}
 \le \left\|\Law(Z_N)-Q_{\mu_N,\alpha_N}\right\|_1
 \le C_{C,\varepsilon}N^{-4/3+\varepsilon}.
 \label{p1:eq:main}
\end{equation}
The proof is unconditional.
Moreover, \(\mu_N-\lambda_N=O_{C,\varepsilon}(N^{-1+\varepsilon})\).
The reference \(Q_{\mu_N,\alpha_N}\) is a probability measure for
all sufficiently large \(N\).
\end{theorem}

The exponent \(4/3\) is optimal with an arbitrary fixed power
margin. The theorem does not determine a logarithmic factor or
a leading constant for this remainder. Its parameters are the
exact mean and second factorial cumulant. Replacing \(\mu_N\)
by \(\lambda_N\) costs \(O(N^{-1+\varepsilon})\), so this
replacement gives a different precision.

The arithmetic calculation in Theorem~\ref{p1:arith:pairs} gives
the first two finite channel contributions to \(\alpha_N\),
with error \(O(N^{-4/5+\varepsilon})\).
Section~\ref{p13:channels:section} extends this calculation
to every fixed height. In particular, if \(r\) denotes the residue
of \(L\) modulo \(3\), then
\begin{equation}
 \alpha_N\sim
  \frac{r+2}{18\,2^{r/3}}\lambda_N^{5/3}L N^{-2/3}.
 \label{p1:eq:alpha-leading}
\end{equation}
The ratio in this equivalent tends to one along the full critical
band, with the displayed moving parameters retained.

\begin{corollary}[Exact total variation scale]
\label{p1:cor:tv}
Under the same assumptions,
\begin{equation}
 \TV(\Law(Z_N),\Pi_{\lambda_N})
   =\frac{\alpha_N}{4}
       \|\Pi_{\lambda_N}*\Delta^{*2}\|_1
       +O_{C,\varepsilon}(N^{-1+\varepsilon}).
 \label{p1:eq:tv-main}
\end{equation}
Consequently,
\begin{equation}
 \TV(\Law(Z_N),\Pi_{\lambda_N})
 \sim
 \frac{r+2}{72\,2^{r/3}}\,
 \lambda_N^{5/3}
 \|\Pi_{\lambda_N}*\Delta^{*2}\|_1\,
 L N^{-2/3}.
 \label{p1:eq:tv-equiv}
\end{equation}
With \(\Pi_{\mu_N}\) as the uncorrected reference, the
first formula holds with \(\mu_N\) in the difference norm
and error \(O_{C,\varepsilon}(N^{-4/3+\varepsilon})\);
the equivalent \eqref{p1:eq:tv-equiv} is unchanged.
\end{corollary}

Here the exponent \(2/3\), the factor \(L\), and the leading
coefficient are exact for the scalar total variation distance
to either of the two displayed Poisson references.
These statements compare with the two specified Poisson means.
They do not minimize total variation over all possible means.

\begin{example}[A numerical coefficient without a moving phase]
\label{p1:ex:unit}
Take \(N=2^{3k}\) and \(L=3k\). Then \(\lambda_N=1\) and
\(r=0\). Since
\[
 \|\Pi_1*\Delta^{*2}\|_1=\frac3e,
\]
Corollary~\ref{p1:cor:tv} becomes
\begin{equation}
 \TV(\Law(Z_N),\Pi_1)
       \sim \frac{1}{12e}\,L N^{-2/3}.
 \label{p1:eq:unit-example}
\end{equation}
The factor \(1/12e\) combines an arithmetic coefficient \(1/9\)
with the second-difference norm \(3/e\) and the factor \(1/4\)
from the cumulant and total variation conventions.
\end{example}

An entirely explicit reference has a different optimal rate.
Retain the complete \(2:3\) and \(3:4\) channel sums, and set
\begin{equation}
 \begin{aligned}
 \widehat\alpha_N
   &=\frac{\lambda_N^2}{3N}\{S_{2,3}(L)+S_{3,4}(L)\},\\
 S_{a,b}(L)
   &=\sum_{d=1}^{\lfloor L/b\rfloor}A_d(L-1-ad)(L-1-bd).
 \end{aligned}
 \label{p1:eq:alpha-explicit}
\end{equation}
where \(A_1=1\) and \(A_d=2^{d-2}\) for \(d\ge2\); an empty
sum is zero. The closed formulas
\eqref{p1:arith:S23}--\eqref{p1:arith:S34} evaluate these sums.

\begin{corollary}[An explicit probability reference]
\label{p1:cor:explicit}
For every \(\varepsilon>0\), uniformly in the critical band,
\begin{equation}
 \left\|\Law(Z_N)-\Pi_{\lambda_N}
       -\frac{\widehat\alpha_N}{2}
          \Pi_{\lambda_N}*\Delta^{*2}\right\|_1
       \ll_{C,\varepsilon}N^{-4/5+\varepsilon}.
 \label{p1:eq:explicit-law}
\end{equation}
For all sufficiently large \(N\), the reference
\(Q_{\lambda_N,\widehat\alpha_N}\) is a probability measure.
Its mass at \(k\ge0\) is
\begin{equation}
 e^{-\lambda_N}\frac{\lambda_N^k}{k!}
 \left[
  1+\frac{\widehat\alpha_N}{2}
       \left\{\frac{k(k-1)}{\lambda_N^2}
                    -\frac{2k}{\lambda_N}+1\right\}
 \right].
 \label{p1:eq:explicit-masses}
\end{equation}
All parameters in this reference are determined by \(N,L\).
\end{corollary}

\begin{proof}
Theorem~\ref{p1:arith:pairs} gives
\(\alpha_N-\widehat\alpha_N=O_{C,\varepsilon}(N^{-4/5+\varepsilon})\).
Since \(\|\Pi_{\lambda_N}*\Delta^{*2}\|_1\le4\), replacing the
coefficient costs at most \(2|\alpha_N-\widehat\alpha_N|\).
Apply Theorem~\ref{p1:thm:main}.
The replacement of \(\mu_N\) by \(\lambda_N\) costs
\(O(N^{-1+\varepsilon})\), which is absorbed here.
The finite formulas give
\(\widehat\alpha_N\sim\alpha_N>0\) and
\(\widehat\alpha_N\to0\), uniformly in the band.
The positivity criterion in Lemma~\ref{p1:close:positivity}
therefore applies, since \(\lambda_N\) remains in a compact
subinterval of \((0,\infty)\).
\end{proof}

The exponent \(4/5\) in \eqref{p1:eq:explicit-law} is optimal
for this two-channel reference. Its leading error is half the
height-five contribution times
\(\|\Pi_{\lambda_N}*\Delta^{*2}\|_1\), as proved in
Corollary~\ref{p13:channels:height-five}.
The explicit parameters do not supply a numerical value of
the constant in the remainder or a numerical threshold for \(N\).
The same two-channel coefficient also gives a positive Hermite
reference; see Corollary~\ref{p1:cor:compound}.

\paragraph{The probability of no retained start.}
Evaluating \eqref{p1:eq:explicit-law} at zero gives the immediate
consequence
\begin{equation}
 \P(Z_N=0)=e^{-\lambda_N}
       \left(1+\frac{\widehat\alpha_N}{2}\right)
       +O_{C,\varepsilon}(N^{-4/5+\varepsilon}),
 \label{p1:eq:no-start}
\end{equation}
because \((\Pi_{\lambda_N}*\Delta^{*2})\{0\}=e^{-\lambda_N}\).
The positive correction shows an effect of the arithmetic pairs
that is absent from the mean alone. Here \(Z_N=0\) means that
no retained run start lies in \(I_N\). A run starting before
\(N\) may still intersect the interval, and a run starting in
\(I_N\) may continue beyond \(2N\).
With exact parameters, the more precise form is
\[
 \P(Z_N=0)=e^{-\mu_N}(1+\alpha_N/2)
                  +O_{C,\varepsilon}(N^{-4/3+\varepsilon}).
\]

\paragraph{Two further optimal powers.}
For a fixed integer \(h\ge3\), retain every reduced interior
channel \(a:b\) with \(b\le h\), and write
\[
 \widehat\alpha_N^{(\le h)}
   =\frac{\lambda_N^2}{N}
      \sum_{b=3}^{h}
      \sum_{\substack{b/2<a<b\\(a,b)=1}}
                  \frac{2(2a-b)}{ab}S_{a,b}(L).
\]
Theorem~\ref{p13:channels:hierarchy} gives an equivalent for
\(\|\Law(Z_N)-Q_{\lambda_N,\widehat\alpha_N^{(\le h)}}\|_1\)
from the first omitted height \(h+1\), of order
\(LN^{-h/(h+1)}\) with an explicit moving coefficient.
Thus the powers \(3/4,4/5,5/6,\ldots\) distinguish finite
explicit references, while \(4/3\) concerns exact cumulants.
The height \(h\) is fixed, not a function of \(N\).

For the third factorial cumulant
\[
 \kappa_{3,N}=\E(Z_N)_3-3\mu_N\E(Z_N)_2+2\mu_N^3,
\]
Theorem~\ref{p13:third:main} proves
\[
 c_C N^{-41/30}\le\kappa_{3,N}
                  \ll_{C,\varepsilon}N^{-41/30+\varepsilon}.
\]
The lower bound comes from genuine compatible triples in
the rational configuration \(10:12:15\). The upper bound
includes all triples, without a conjecture on radicals or
smooth numbers. It does not identify a full asymptotic
constant for this cumulant.

\begin{table}[!htb]
\centering
\small
\caption{References and established precision in the critical band.
Power bounds allow an arbitrary fixed \(\varepsilon>0\).}
\label{p16:tab:references}
\begin{tabularx}{\textwidth}{@{}>{\raggedright\arraybackslash}p{0.25\textwidth}>{\raggedright\arraybackslash}p{0.20\textwidth}>{\raggedright\arraybackslash}X@{}}
\toprule
Reference or observation & Parameters & Established precision\\
\midrule
Uncorrected Poisson
& \(\lambda_N\) or \(\mu_N\)
& Explicit equivalent of order \(LN^{-2/3}\)
(Corollary~\ref{p1:cor:tv}).\\[4pt]
Exact quadratic \(Q\)
& \(\mu_N,\alpha_N\)
& Optimal power \(4/3\); no full remainder equivalent
(Theorem~\ref{p1:thm:main}).\\[4pt]
Nominal-mean quadratic \(Q\)
& \(\lambda_N,\alpha_N\)
& Upper bound \(O(N^{-1+\varepsilon})\); comparable to
\(|\mu_N-\lambda_N|\) when this bias is nonzero
(Proposition~\ref{p13:optimal:nominal}).\\[4pt]
Quadratic with channels through height \(h\)
& \(\lambda_N,\widehat\alpha_N^{(\le h)}\)
& Equivalent of order \(LN^{-h/(h+1)}\) from the first
omitted height, for each fixed \(h\ge3\)
(Theorem~\ref{p13:channels:hierarchy}).\\[4pt]
Exact Hermite
& \(\mu_N,\alpha_N\)
& Upper bound \(O(N^{-4/3+\varepsilon})\); no matching
optimal lower bound established here
(Corollary~\ref{p1:cor:compound}).\\[4pt]
Unsigned positional field
& Individual site and edge intensities
& Upper bound \(O(N^{-4/5+\varepsilon})\) for each fixed
polynomial weight in the total count; signs and excess lengths are not retained
(Lemma~\ref{p1:limits:vector-reference}).\\
\bottomrule
\end{tabularx}
\end{table}
\FloatBarrier

\subsection{Why a pair has the leading size}
\label{p1:intro:mechanism}

The dyadic geometry selects the first interior channel.
If \(x\in[N,2N)\), then \(2x\notin[N,2N)\), so the
ratio \(1:2\) has no macroscopic family of interior placements;
offsets of size \(O(L)\) leave only a boundary contribution.
The first interior channel is \(2:3\), as the general interval
calculation in Proposition~\ref{p1:consequences:macro} confirms.

Consider \(L=3m\), and centres \(x=2t\), \(y=3t\) in \(I_N\).
For \(0\le j\le m-1\), the values \(2(t+j)\) and \(3(t+j)\)
belong to their respective windows. For \(j\ne k\),
\[
 f(2(t+j))f(2(t+k))
       =f(3(t+j))f(3(t+k)).
\]
There are \(m-1\) independent relations of this form. When these
are the only relations of the pair, its joint probability is
\[
 \P(J_{2t}=J_{3t}=1)=2^{-2L+m-1},
\]
rather than \(2^{-2L}\). The preceding signs do not occur in these
relations, so the two run-start prescriptions are compatible.
The arithmetic section proves that additional relations form a
smaller exceptional family.

Allowing the offset \(2y-3x\) to range over its admissible interval
of length \(O(L)\) produces \(O(NL)\) pairs. A typical joint weight
on this channel is \(N^{-5/3}\). Their total weight is therefore
of order \(L N^{-2/3}\). The analogous \(3:4\) channel has joint
weight \(N^{-7/4}\) and total weight \(L N^{-3/4}\).
Summing the exact endpoint and parity conditions yields
\eqref{p1:eq:alpha-leading} and the finite coefficient in
Theorem~\ref{p1:arith:pairs}.

The full finite contribution of \(2:3\) matters even at the precision
of Corollary~\ref{p1:cor:explicit}: its next term, of order \(N^{-2/3}\),
is larger than the leading \(3:4\) term. Thus the second channel
is isolated by subtracting the entire \(2:3\) expression.
The two finite expressions have a combined phase modulo \(12\);
only the leading equivalent has phase modulo \(3\).

\paragraph{The ternary witness and the residual quadratic error.}
For the centres \(10t,12t,15t\), multiplication by \(6,5,4\),
respectively, gives the common centre \(60t\) and offset steps
\(6,5,4\). Their rational nullity is
\[
 \rho_L=\left\lfloor\frac{L-1}{6}\right\rfloor
       +\left\lfloor\frac{L-1}{3}\right\rfloor
       +\left\lfloor\frac{L-1}{5}\right\rfloor
       -\left\lfloor\frac{L-1}{15}\right\rfloor
       =\frac{19}{30}L+O(1).
\]
Thus a compatible triple with no additional relations has weight
of order \(N^{-3+19/30}=N^{-71/30}\); a family of order \(N\)
translations explains the power \(41/30\).
Proposition~\ref{p13:third:lower} makes this witness rigorous by
checking private coordinates, predecessor signs, positivity of the
centered coefficient, and exceptional hosts. The upper estimate for
all triples is a separate part of Section~\ref{p13:third:section}.
In particular, the coefficient \(1/120\) obtained when \(N=2^L\)
and \(30\mid L\) describes the sum of centered coefficients over
the selected translation subfamily, not the whole cumulant.

Two pair interactions have the larger scale
\(\alpha_N^2\asymp L^2N^{-4/3}\). To see where their fourth-order
term enters, compare the quadratic reference with the Hermite law
\[
 H_{\mu,\alpha}=\Pi_\mu*\exp_*(\alpha\Delta^{*2}/2)
  =Q_{\mu,\alpha}
   +\frac{\alpha^2}{8}\Pi_\mu*\Delta^{*4}
   +O_{\ell^1}(\alpha^3).
\]
This expansion locates a possible obstruction but does not prove
that it survives in the true law. Section~\ref{p13:lower:section}
establishes the noncancellation identity
\[
 \kappa_{4,N}+3\alpha_N^2
   =24K_4^++3\alpha_N^2+O(N^{-3/2+\varepsilon}),
 \qquad K_4^+\ge0,
\]
where \(K_4^+\) sums the centered coefficients over compatible
connected quadruples. Independently bounded fixed moments turn this
unbounded Charlier test into the lower norm bound.
The resulting conclusion is the optimal power \(4/3\),
with an arbitrary fixed exponent margin, rather than an
equivalent for the quadratic remainder.

\subsection{From arithmetic counts to a distance between laws}
\label{p1:intro:strategy}

The proof begins with the exact marginal probabilities
\(p_x=\E J_x\) on all admissible windows. Local rigidity gives
\(p_x\in\{p,2p\}\); an impossible event is omitted without
changing \(Z_N\). An occurrence is an affine condition on the
prime signs over \(\F_2\). Dependencies are measured between
the resulting effective spaces, so local relations are not
counted twice. Independent Bernoulli variables with these
marginals supply the reference; repeated uses of the same
event are thereby retained exactly.

The reconstruction has three operations: retain families with a
calculable joint probability, stop at the first window that leaves
this class, and count the first excluded configurations.
Compression treats a dependent pair as one affine block.
The retained families form a hereditary class. Developing
\(\delta_{J_x}=\delta_0+J_x\Delta\) only until the first
excluded prefix gives the exact remainder
\[
 \sum_{I\in\partial\mathscr D}\Delta^{*|I|}*\nu_I,
 \qquad \nu_I\ge0,\qquad \|\nu_I\|_1=q(I).
\]
All indicators after that prefix stay inside the positive
measure \(\nu_I\), with their dependence on the prefix intact.
Example~\ref{p1:stop:example} makes this point explicit.
Signed pair terms are subtracted before absolute values are taken.

Theorem~\ref{p13:transfer:reconstruction} proves
\[
 \|\Law Z_N-Q_{\mu_N,\alpha_N}\|_1
 \ll_{C,\varepsilon}N^\varepsilon A_3+N^{-4/3+\varepsilon},
 \qquad
 A_3=\sum_{|T|=3}\left|\E\prod_{x\in T}(J_x-p_x)\right|.
\]
Its finite-order estimates and its uniform tail are proved
separately. Small primitive configurations are reduced to
triples. For large configurations, moderate divisibility
certificates are counted by congruences; oversized certificates
have small incoming and outgoing degrees. A component-generating
function then sums the tail. A bound for each fixed order alone
would not justify this step.

The arithmetic problem is now to bound \(A_3\) strictly below
power \(4/3\). The conditional half-power count in
Theorem~\ref{p13:half:conditional} makes incompatible triple
activity \(O(N^{-3/2+\varepsilon})\). Compatible triple
coefficients are nonnegative. Their partition by pair ranks
reduces the estimate to rational alignments and rare pairs
without such an alignment.

A nonaligned pair cannot share half of a window's constraints.
Short relations extracted from its prime-factor equations force
a common algebraic curve through the two centres. The relations
also force many integer points in the much smaller offset box.
The degree-uniform bound of Binyamini, Cluckers, and Novikov
\cite[Theorem~2, (27)]{binyaminicluckersnovikov2024} gives
\[
 \dim(S_x\cap S_y)\le
       \left(\tfrac23\log2+o(1)\right)L<\tfrac12L.
\]
Rational triple geometry then yields
\(A_3\ll N^{-41/30+\varepsilon}\).
The gain \(41/30-4/3=1/30\) closes the upper bound after
choosing a sufficiently small fixed error margin.
The compatible \(10:12:15\) family proves that \(41/30\)
is itself the optimal power of the third factorial cumulant;
this lower bound is a separate conclusion, not an input to
the upper approximation theorem.

Finally, a fourth Charlier polynomial tests the residual law
and detects \(3\alpha_N^2\), up to a smaller error and a
nonnegative compatible contribution. Since
\(\alpha_N^2=N^{-4/3+o(1)}\), interpolation with independently
bounded fixed moments proves the lower power in
Theorem~\ref{p1:thm:main}. Both bounds concern the same
reference with exact parameters.

The stopping identity applies to any finite family of indicators.
Its use in another model requires a retained class with a
calculable reference and a small boundary, together with
uniform estimates for that boundary. Here these estimates
come from the arithmetic geometry of the windows.

\subsection{Context and contribution}

The preceding work on this multiplicative model proves a Poisson
approximation for the entire sign-resolved run field
\cite[Theorem~1.1]{pouly2026c}. We study the scalar count
\eqref{p1:eq:model} at a finer scale. Summing the two signs removes
the leading one-value square bias; the first surviving correction
comes from relative relations between pairs of windows.
The sharp coefficient in \eqref{p1:eq:tv-equiv} concerns the
scalar count. Lemma~\ref{p1:limits:vector-reference} also compares
the unsigned field \((J_x)_{x\in I_N}\) on its integer grid with
independent Poisson single arrivals and arrivals along good
arithmetic pairs. Its error is
\(O_{C,j,\varepsilon}(N^{-4/5+\varepsilon})\) in \(\ell^1\)
weighted by \((1+\text{total count})^j\), for every fixed
\(j\ge0\). This comparison retains positions and justifies
deletion of doublets; it does not retain run signs or exact
excess lengths. The target uses individual edge weights
\(q_e=\P(J_x=J_y=1)\), \(e=\{x,y\}\), which carry information beyond the
scalar coefficient \(\widehat\alpha_N\).
No corresponding sharp distance equivalent is asserted for the
sign-resolved field or for a diffuse spatial point process.

The interaction \(2:3\) already occurs in the lower-bound
argument of that work \cite[Section~3.6, equation~(3.22)]{pouly2026c}.
Here we compute its complete contribution after summing the
offsets and affine signs, retain the secondary channel \(3:4\),
and control the remaining law below the size of the leading
correction. Detecting a correlated pair is only the first part
of that calculation.

For independent sequences, long runs and rare patterns have a
classical Poisson theory; see Arratia, Gordon, and Waterman
\cite{agw1990} and Godbole \cite{godbole1991}. Counting starts
separates occurrences that belong to one long run. In the present
model, multiplicative relations also couple separated windows.
Najnudel \cite[Lemma~3.1 and Proposition~6.3]{najnudel2020}
describes consecutive values through valuation relations and
proves almost-sure frequencies of fixed words for random
completely multiplicative functions with root-of-unity values
at primes. The fixed-length frequency statement does not
provide the uniform rare-event estimates used here when
\(L\asymp\log N\). In the binary model, independence of coordinates
also has to be distinguished from fairness of every coordinate:
values at squares are deterministic.

A related arithmetic Poisson phenomenon is studied by Basak,
Nath, and Zaharescu \cite[Theorem~1]{basaknathzaharescu2024}.
They count blocks of consecutive quadratic non-residues modulo
a prime, and obtain Poisson frequencies under translation for
almost all primes in suitable intervals. Their blocks lie in
an initial polylogarithmic range, with a length tending to
infinity no faster than four iterated logarithms of the prime
scale. Their source of averaging and their block count differ
from the prime-sign probability space and run-start count
\eqref{p1:eq:model}. In another direction, de la Bret\`eche,
Wang, and Xu \cite{bretwngxu2026} prove central limit theorems
for sums of random multiplicative functions on polynomial
values by counting square products. Their completely
multiplicative binary model is called extended Rademacher;
the squarefree-supported Rademacher convention is a different
model. The shared square-relation tools serve different
observables in those central limit theorems and in the
run-counting problem.

Poisson approximation for dependent indicators has a general
Chen--Stein framework \cite{agg1989}. Refinements of Poisson
limits include Barbour's asymptotic expansions
\cite{barbour1987} and the expansions for dissociated variables
of Barbour and Eagleson \cite{barboureagleson1987}.
The Poisson--Charlier form and its organization by factorial
cumulants are classical; see Barbour, Kowalski, and Nikeghbali
\cite[Sections~5 and~8.1]{bkn2014}, in the author version
specified in the bibliography. The difference norms used below
are recorded by Zacharovas and Hwang
\cite[equation~(2.13) and Remark~4.4]{zacharovashwang2010},
and the single- and double-jump reference is the Hermite law
of Kemp and Kemp \cite{kempkemp1965}.

\v{C}ekanavi\v{c}ius and Vellaisamy
\cite[Theorems~3.1--3.2]{cekanaviciusvellaisamy2015} compare
quadratic and exponential corrections for locally dependent sums;
their remainder estimates exhibit the cost of truncating the
second-cumulant exponential. Their local-dependence hypotheses
do not cover the arithmetic interactions between separated windows
used here. Chhaibi, Delbaen, M\'eliot, and Nikeghbali
\cite[Theorem~3.11]{chhaibietal2020} obtain distance equivalents
for discrete approximations in a mod-\(\phi\) framework. That
theorem assumes a parameter tending to infinity and a regularity
condition on the Fourier residues. Here the Poisson means remain in a compact
subset of \((0,\infty)\), and the error is controlled by direct
arithmetic estimates.

Our contribution is the determination of the leading arithmetic
deviation from Poisson approximation for the scalar run-start
count, the optimal quadratic remainder power, and the
hierarchy of errors for explicit channel truncations.
The third-cumulant exponent and the pair rigidity estimate
identify which configurations can contribute at those scales.
The finite stopping identity is a specialization of the abstract
lace expansion \cite{zeilberger1997}, as detailed in
Section~\ref{p1:stop:section}.
Its role here depends on a retained class adapted to the pair
geometry and on uniform estimates for the boundary of that class.
Those estimates, together with the arithmetic coefficient, are
the model-specific work.

\subsection{Proof dependencies and ranges of uniformity}
\label{p1:intro:dependencies}

This article and its \emph{Technical companion} form one work.
Numbered references beginning with T refer to the companion;
other references refer to this article. The companion supplies
full technical proofs of the inputs stated here and the further
results on tails, deletion, and conditional geometry. Both
documents are needed to check the complete proof. The article
follows the exact-marginal argument directly.

Throughout, the critical-band constant \(C\) and every requested
error margin \(\varepsilon_0>0\) are fixed before \(N\to\infty\).
The following map separates the inputs to the upper and lower bounds.
Each arrow denotes the use of the entries in its own row.

\begin{center}
\small
\begin{tabularx}{\textwidth}{@{}>{\raggedright\arraybackslash}p{0.41\textwidth}>{\centering\arraybackslash}p{0.035\textwidth}>{\raggedright\arraybackslash}X@{}}
\toprule
Inputs & & Conclusion and uniformity\\
\midrule
Fixed-source arithmetic and propagation
(Sections~\ref{p1:arith:appendix}, \ref{p14:comp:fixed})
& \(\longrightarrow\)
& Fixed moments, Theorem~\ref{p1:fixed:moments}.
The order is fixed; the proof precedes every law approximation.\\[5pt]
Short-selection moments, conditional half counts, pair rigidity
(Sections~\ref{p13:short:section}--\ref{p13:rig:section})
& \(\longrightarrow\)
& \(A_3\ll N^{-41/30+\eta}\),
Corollary~\ref{p13:third:activity}.
All admissible triples are included.\\[5pt]
Triple activity; fixed connected orders through \(256\);
separate certificate tails
(Section~\ref{p13:transfer:section})
& \(\longrightarrow\)
& Full-law reconstruction,
Theorem~\ref{p13:transfer:reconstruction}, and the upper bound
in Theorem~\ref{p1:thm:main}.\\[5pt]
Pair profile and fixed-source counts, triple activity,
conditional extension deficits
(Lemmas~\ref{p13:transfer:independent-deficit}
and~\ref{p13:transfer:fixed-deficit})
& \(\longrightarrow\)
& The noncancelling quartic identity,
Lemma~\ref{p13:lower:quartic}.
It uses these counting estimates, not the reconstruction conclusion.\\[5pt]
Quartic identity and the independently bounded fixed moments
& \(\longrightarrow\)
& The lower bound, Theorem~\ref{p13:lower:optimal}.
The interpolation order is chosen after \(\varepsilon_0\)
and remains fixed.\\[5pt]
Main theorem and the same fixed moments
& \(\longrightarrow\)
& Polynomially weighted scalar errors,
Corollary~\ref{p13:optimal:weighted}; these are consequences.\\
\bottomrule
\end{tabularx}
\end{center}

The moments of short selections in Section~\ref{p13:short:section}
concern multiplicative functions on selected polynomial values.
Their degree may reach \(c\log\log N\). They are distinct from the
fixed moments of \(Z_N\) in Theorem~\ref{p1:fixed:moments}.
The curve degree in Section~\ref{p13:rig:section} has its own
controlled growing range. Neither estimate is obtained by allowing
the order in a fixed-moment theorem to grow.

A separate branch retains the regular-set reconstruction in
Sections~\ref{p14:comp:stopped}--\ref{p1:close:section} of the
companion. Its weighted estimates give the unsigned positional
comparison in Section~\ref{p1:limits:positions} and
Lemma~\ref{p1:limits:vector-reference}, with power \(4/5\).
The channel hierarchy in Section~\ref{p13:channels:section}
also keeps its own quantifier: the height \(h\) is fixed.

The quasi-linear pair bound in Proposition~\ref{p1:arith:hosts} is
proved only after the conditional half-count and confinement; its
earlier elementary $6/5$ version suffices for the pair-profile branch.
The fixed-arity anchored argument also gives the stronger primitive
weights in Proposition~\ref{p13:transfer:primitive}. Its parameters
$r$ and, in the dense sector, $\eta$ are fixed. These improvements
do not replace the separate all-cardinality tail or remove the
two-pair term at the $4/3$ scale.

The cutoff \(256\) is sufficient; no optimality is claimed for it.
The arithmetic inputs in Sections~\ref{p1:arith:appendix}
and~\ref{p1:arith:supplement} are part of both reconstruction branches.

\paragraph{Power margins.}
The strict power savings used in the last stages are recorded below.
Each margin compares complete contributions, after coefficient choices,
propagation and rank factors have been included.
\begin{center}
\small
\begin{tabularx}{\textwidth}{@{}Xcl@{}}
\toprule
Comparison & Exact margin & Location\\
\midrule
Nonaligned target above a rational pair versus dominant triples
& $11/8-41/30=1/120$
& \eqref{p13:third:unaligned-target-cost}\\[4pt]
Incompatible triples versus the third factorial cumulant
& $3/2-41/30=2/15$
& \eqref{p13:activity:kappa-comparison}\\[4pt]
Third activity versus the quadratic remainder
& $41/30-4/3=1/30$
& \eqref{p13:transfer:reconstruction-bound}\\[4pt]
Low-nullity uniform tail versus the quadratic remainder
& $27/20-4/3=1/60$
& \eqref{p13:transfer:low-rank-tail}\\[4pt]
Quartic signed error versus $\alpha_N^2$
& $3/2-4/3=1/6$
& \eqref{p13:lower:positive-quartic}\\
\bottomrule
\end{tabularx}
\end{center}
The smallest displayed margin is $1/120$. In contrast, two pair
activities have size $A_2^2\ll N^{-4/3+\eta}$ and are retained at
the target scale. They have no strict power saving to spend.
For a prescribed final accuracy $\varepsilon_0>0$, one may choose
$0<\tau<\min(\varepsilon_0/4,1/960)$ and invoke completed auxiliary
estimates with total loss at most $\tau$ each. Their own input margins
are chosen smaller as needed; in a product, divide the allotted loss
between the factors. Thus the same symbol $\varepsilon$ in successive
statements never requires the same numerical value at every stage.
Rank thresholds and fixed geometric degrees have separate roles.
The fixed connected-activity estimates in reconstruction are used only
through order $256$; larger cardinalities have the uniform certificate estimate and a factorial
boundary bound.

\subsection{Notation guide}
\label{p1:intro:notation}

The following guide distinguishes the sets of integer arguments
from their spaces of square classes. The full definitions and
the compatibility conditions are given where each object is used.

\begin{center}
\small
\begin{tabularx}{\textwidth}{@{}>{\raggedright\arraybackslash}p{0.29\textwidth}>{\raggedright\arraybackslash}X@{}}
\toprule
Notation & Meaning\\
\midrule
\(\lambda_N,\mu_N,\alpha_N\)
& \(N2^{-L}\), \(\E Z_N\), and
\(\Var(Z_N)-\mu_N\), respectively.\\[3pt]
\(\widehat\alpha_N\)
& The explicit two-channel coefficient in
\eqref{p1:eq:alpha-explicit}, whose first omitted height
is five.\\[3pt]
\(\beta_N=\sum_{e\in E_N}q_e\)
& The positive mass of good pairs in the regular-set branch.
Its conversion to half
the factorial cumulant uses the covariance and Bernoulli
diagonal corrections in Section~\ref{p1:close:section}.\\[3pt]
\((z)_r\)
& The falling factorial, with \((z)_0=1\).\\[3pt]
\(\rho_{\mathrm{rel}}(I),p_x\)
& The effective relative rank
\(\sum_{x\in I}\dim S_x-\dim\sum_{x\in I}S_x\)
and the exact marginal \(\E J_x\).
They are used from the outset; on the regular set they
reduce to the homogeneous quantities.\\[3pt]
\(V_x\), also \(\mathcal V_x\)
& The set \(\{x-1,\ldots,x+L-1\}\) of integer
arguments in a window.\\[3pt]
\(E_x,S_x\)
& The space of even selections of \(V_x\), and its image
under prime valuations modulo \(2\).
On \(G_N\), \(\dim S_x=L\).\\[3pt]
\(\rho(I)\), \(\rho_*(I)\)
& In the pair arithmetic, \(\rho(I)=|I|L-\dim\sum_{x\in I}S_x\),
including local kernels. It equals \(\rho_{\mathrm{rel}}\) on \(G_N\).
The compatible nullity \(\rho_*\) is the maximum effective nullity
of a compatible subfamily. Section~\ref{p13:transfer:section}
explicitly abbreviates \(\rho_{\mathrm{rel}}\) to \(\rho\).\\[3pt]
\(q(I)\), \(w(I)\), \(W(I)\)
& \(q(I)=\eta(I)w(I)\), where \(\eta(I)\in\{0,1\}\)
indicates compatibility and
\(w(I)=2^{\rho_{\mathrm{rel}}(I)}\prod_{x\in I}p_x\).
The regular-set companion uses \(W(I)=p^{|I|}2^{\rho(I)}\).\\[3pt]
\(c(I)\)
& \(\E\prod_{x\in I}(J_x-p_x)\); on \(G_N\),
all \(p_x=p\).\\[3pt]
\(U_e,V_e\)
& For a pair \(e=\{x,y\}\) in channel \(a:b\),
\(U_e=b\mathcal V_x\cup a\mathcal V_y\) is a set
of integer values; \(V_e=S_x+S_y\) is a vector space.\\[3pt]
\(\mathscr D,\partial\mathscr D\)
& The retained hereditary class and its boundary of
first excluded prefixes for the chosen ordering.\\[3pt]
\(n,\mu\)
& In the regular-set argument, \(n=|G_N|\) and
\(\mu=np\); this \(\mu\) is distinct from \(\mu_N\).\\
\bottomrule
\end{tabularx}
\end{center}

\paragraph{Reading routes.}
For the statements and their mechanism, read this introduction,
the complete pair profile in Theorem~\ref{p1:arith:pairs},
the reconstruction and lower-bound theorems in
Sections~\ref{p13:transfer:section} and~\ref{p13:lower:section},
and the consequences in Sections~\ref{p13:optimal:section}
and~\ref{p13:channels:section}.
For the proof, follow the article in order and use each explicitly
identified companion proof when its technical input is needed.
The companion's opening guide separates those inputs from the
regular-set reconstruction and the further consequences.

